\documentclass{article} % For LaTeX2e
\usepackage{iclr2026_conference,times}

% Optional math commands from https://github.com/goodfeli/dlbook_notation.
\input{math_commands.tex}

\usepackage{microtype}
\usepackage{graphicx}
\usepackage{amsmath,amssymb}
\usepackage{booktabs}
\usepackage{hyperref}
\usepackage{url}
\usepackage{xspace}

\newcommand{\invarlock}{\textsc{InvarLock}\xspace}
\newcommand{\schemaverify}{\texttt{schema-verify}\xspace}

\title{A Machine-Readable Schema for Weight-Edit Evaluation Reports\\with First-Class Verifier Traces}

% Authors must not appear in the submitted version.
\author{Anonymous Authors}

%\iclrfinalcopy % Uncomment for camera-ready version, but NOT for submission.
\begin{document}

\maketitle

\begin{abstract}
Evaluations of weight edits such as quantization and pruning increasingly mix
statistical guard frameworks (e.g., \invarlock), general evaluation harnesses,
and task-specific external verifiers (e.g., HumanEval/MBPP execution).
Today these tools emit incompatible artifacts, making it difficult to align guard
evidence with verifier outcomes, reproduce end-to-end verification loops, or
build automated promotion gates with auditable inputs.
We propose a lightweight JSON Schema for \emph{verifier-carrying evaluation artifacts}
that links three evidence elements in a single machine-checkable object:
\texttt{guard\_evidence} (tool-native guard measurements, thresholds, and verdicts,
including \invarlock evaluation reports and optional \texttt{invarlock verify} outputs),
\texttt{verifier\_traces} (external verifier results together with a trace contract
covering prompt hashes, harness/sandbox identity, resource limits, and decoding
parameters), and \texttt{provenance} (model/tokenizer identity, edit configuration, and
tool versioning), under a versioned wrapper (\texttt{schema\_version}).
A companion \schemaverify utility validates structural conformance and cross-block
consistency, enabling downstream automation to treat the artifact as a reliable unit
of evidence in verification-aware pipelines.
We implement serializers for \invarlock artifacts and code-execution verifiers and
demonstrate cross-tool queries on a paired dataset spanning multiple models and edit
families, including cases where guard evidence and verifier outcomes disagree.
We release the schema, serializers, validator, and a queryable demo corpus as open-source
building blocks for composing soft evidence with hard verification in model-edit workflows.
\end{abstract}

\section{Introduction}
When a language model is quantized, pruned, or otherwise weight-edited, evaluation evidence is
produced at three layers that rarely interoperate:
(i) \textbf{guard frameworks} (e.g., \invarlock) emit diagnostic measurements and gate flags,
(ii) \textbf{evaluation harnesses} define prompt sets and generation parameters,
and (iii) \textbf{external verifiers} provide the decisive outcome signal (did outputs satisfy the
verifier?).
This structure aligns with verifier-first evaluation: the external verifier outcome is the
headline metric, while guard evidence is auxiliary and diagnostic.
In practice, however, the layers emit incompatible artifacts with missing provenance, and basic
comparability questions become hard to answer automatically:
are two reported pass@1 values computed over the same prompt set, under the same harness revision,
with the same decoding parameters, and under the same execution sandbox constraints?

Existing formats do not close this gap.
Model cards and leaderboards rarely encode verifier determinism contracts \citep{mitchell2019model_cards}.
Dataset metadata standards improve dataset provenance but do not pin verifier-run provenance
\citep{akhtar2024croissant}.
Supply-chain frameworks provide content-addressing primitives that inspire our design but target
software artifacts rather than model-evaluation evidence \citep{slsa2024spec}.

We address this with a \textbf{verifier-carrying artifact}: a strict JSON Schema that co-locates
hash-linked guard evidence with first-class verifier traces under shared provenance.
The verifier trace's \emph{contract} pins comparability-critical conditions (prompt-set digests,
harness identity, model/tokenizer revisions, decoding parameters, and sandbox limits).
\schemaverify checks JSON Schema conformance and contract-driven cross-block consistency, without
claiming formal guarantees about model correctness or verifier soundness.

\textbf{Contributions.}
(1) A verifier-trace schema and contract for auditable external verification runs.
(2) A wrapper schema that packages auxiliary guard evidence and first-class verifier traces in one
queryable object.
(3) A validator that enforces structure and cross-block consistency.

\section{Schema Overview}
The verifier-carrying artifact is a single JSON object
(\texttt{schema\_version="verifier\_carrying\_artifact.v1"}) intended to be the unit of exchange for
weight-edit evaluation evidence.
It is intentionally strict (\texttt{additionalProperties=false}) on core objects so typos and
logging drift fail validation instead of silently passing.
The wrapper has three evidence blocks under shared provenance:
\texttt{provenance}, \texttt{guard\_evidence}, and \texttt{verifier\_traces[]} (minItems=1).
In v1, \texttt{guard\_evidence} includes \texttt{guard\_evidence.invarlock.evaluation\_report} as a
content-addressed file reference (SHA-256) and optional \texttt{invarlock verify} output.
The verifier trace is tool-agnostic: future versions can add peer guard-evidence blocks without
changing the trace contract.

\section{Verifier Trace Contract}
The verifier trace contract is a minimal, machine-checkable record binding an external verifier
outcome (e.g., HumanEval pass@1 \citep{chen2021humaneval} or MBPP pass@1 \citep{austin2021mbpp}) to
the concrete conditions under which it was produced.
It records: verifier identity (name/kind + harness identity), prompt-set integrity (dataset
identifiers plus ordered prompt ids and hashes with a dataset-level digest computed over canonical
JSON bytes), model/tokenizer id+revision, and decoding parameters (including seed).
For code-execution verifiers, the schema requires an explicit sandbox block (network disabled,
timeouts, and resource limits), making execution constraints auditable rather than implicit.

The contract is a trust boundary: it makes comparability and auditability machine-checkable.
It does not ``prove'' determinism across platforms; rather, it defines the minimum information
needed to diagnose divergences and support replay.

\section{Validator}
\schemaverify performs structural validation against Draft 2020-12 JSON Schema for both the wrapper
artifact and each verifier trace.
It also checks contract-defined cross-block invariants, including prompt-set digest recomputation,
alignment between prompt-set ids and per-case result ids, and (optionally) on-disk SHA-256
verification when local paths are available.
These checks establish evidence integrity and joinability, not behavioral correctness.

\section{Demonstration Corpus}
We include a queryable demonstration summary
(\texttt{research/verifai2\_2026/examples/s1\_demo.summary.jsonl}) illustrating joins and audit
queries over verifier outcomes and guard evidence.
The demo contains 11 verifier-carrying artifacts, each paired with two verifier traces (HumanEval
and MBPP), yielding 22 joined rows spanning three base model families (Phi-2, DeepSeek-Coder-1.3B,
and Qwen2.5-Coder-3B) and multiple edit families (INT8 RTN recipes and magnitude pruning masks).

Table~\ref{tab:demo} shows representative rows.

\begin{table}[t]
\centering
\small
\begin{tabular}{llllrrrr}
\toprule
Model family & Edit & Verifier & $\Delta$pass@1 & PM & Spec & RMT & Spec status \\
\midrule
DeepSeek-Coder-1.3B & INT8 RTN (mild) & HumanEval & +0.012 & \checkmark & \checkmark & \checkmark & capped \\
DeepSeek-Coder-1.3B & INT8 RTN (mild) & MBPP     & $-$0.006 & \checkmark & \checkmark & \checkmark & capped \\
Qwen2.5-Coder-3B    & INT8 RTN (mild) & HumanEval & $-$0.006 & \checkmark & \checkmark & \checkmark & capped \\
Phi-2 (BYOE)        & pruning 40\%    & HumanEval & +0.220 & $\times$ & \checkmark & \checkmark & stable \\
\bottomrule
\end{tabular}
\caption{Representative demo rows (joined view). ``Spec status=capped'' means caps were applied
(\texttt{caps\_applied>0}); it is not itself a failure. Gate outcomes are recorded separately.}
\label{tab:demo}
\end{table}

\paragraph{Example queries.}
The joined view enables simple queries that are hard to answer from per-tool logs.
For example: ``primary-metric gate fails, yet HumanEval improves materially'' can be expressed as a
filter over joined records, surfacing guard--verifier disagreement cases that motivate
trust-calibrated escalation policies (consult the verifier before rejecting an edit based solely on
guard gates).

\section{Limitations}
This work standardizes evidence, not guarantees.
The validator checks structural well-formedness and contract-driven consistency; it cannot certify
verifier soundness or harness correctness without rerunning the pipeline.
The demonstration corpus is illustrative; larger-scale analyses of guard--verifier alignment are
out of scope here.

\bibliographystyle{iclr2026_conference}
\bibliography{references}

\end{document}
